\scanpage{44}
\deleted{$\mathfrak X$ \inserted{\unsure{pré}}schéma formel loc. noeth.}

\deleted{$\mathfrak X_0$ un sous-schéma de définition, $X_0=\mathrm V(\unsure{\mathcal J|\mathcal I})$}

\deleted{On \unsure{a|va} définit la catégorie $\mathrm{Coh}(\mathfrak X/\mathfrak X_0)$ des ``\deleted{revêtements finis}\inserted{Modules \unsure{Cohérents|cohérents}} de $\mathfrak X$ mod $X_0$''\unsure{\}}. Soit $\mathrm{Coh}(\mathfrak X)$ la catégorie des \deleted{schémas formels finis}\inserted{modules Coh.} sur $\mathfrak X$. \deleted{Soit} Soit $\unsure{\mathscr F|\mathscr S|\mathfrak S}$ l'ensemble des flèches de $\mathrm{Coh}(\mathfrak X)$ telles que $\mathrm{Ker}\,u$ et $\mathrm{Coker}\,u$ soit \deleted{\unsure{p. à p. par}} annulé \deleted{par}\inserted{loc.} par une puissance de $\mathcal J$. Alors l'ens. $\mathscr F$ \deleted{\unsure{est}} satisfait aux axiomes qu'il faut pour passer à une catégorie de fractions [N.B. on pense en particulier \unsure{évid.} à la sous-catégorie épaisse des Modules coh. tués par \unsure{un|une} $\unsure{\mathcal J^n|\mathcal J}$].}

\deleted{$\mathrm{Coh}(\mathfrak X/\mathfrak X_0)$ est la catégorie des fractions.}

\deleted{\uline{Exemple} Soit $X$ propre sur $Y=\mathrm{Spec}\,A$, $A$ noeth séparé et complet pour idéal $I$, $Y_0=\mathrm V(I)$, $X_0=X\times_Y Y_0$, $\mathfrak X$ le complété formel de $X$.}

\deleted{Alors}

\diagram{\begin{tikzcd}[ampersand replacement=\&,row sep=large]
\deleted{\mathrm{Coh}(\mathfrak X)} \& \deleted{\mathrm{Coh}(X)} \arrow[l,"\simeq"'] \& {\text{\deleted{par le th. d'existence}}}\\
\deleted{\mathrm{Coh}(\mathfrak X)_{\text{\unsure{négl|nilp}}}} \arrow[u,hook] \& \substack{\deleted{\mathrm{Coh}(X)_{\text{négl}}}\\\text{\deleted{(support $\subset X_0$}}} \arrow[u,hook] \&
\end{tikzcd}}

On trouve $\mathrm{Coh}(\mathfrak X/\mathfrak X_0)\cong\mathrm{Coh}(X-X_0)$ !
% archive sheet number: 43
